\chapter{What we cannot claim: limitations and honest answers}
\label{ch:limits}

An evaluator's best questions are about limits. If you have already said the limits yourself, those questions lose their sting. This part is the list, with a one-sentence reply for each.

\section{The limitation list, in plain words}

\begin{enumerate}
\item \textbf{We only see eight numbers per service.} A brownout on a nearly idle service moves none of them enough to alarm. No engine here catches it. \emph{Reply:} ``Outside-in observation is coarse by design; that is the trade for touching nothing. Active probing is listed as future work.''

\item \textbf{Two real systems.} We argue the method is architecture-independent from a cross-system check and from synthetic topologies, but we have not proved it. \emph{Reply:} ``Two systems we built, 217 minutes each, plus 16 random topologies. It is evidence, not proof.''

\item \textbf{Real cascades were rare.} Learned propagation could not be tested on real data because real faults are mostly absorbed single-service outages. \emph{Reply:} ``On synthetic cascades it is better calibrated. On real data there were too few cascades to say.''

\item \textbf{Edge evidence is not isolated offline.} The recordings do not keep it. \emph{Reply:} ``We use the evidence weighting live; we did not measure its isolated effect.''

\item \textbf{Cost was measured on synthetic readings.} Random graphs of 10--320 services on a shared virtual machine. \emph{Reply:} ``About 0.37 ms per service per window, linear; we have not observed a real system of that size.''

\item \textbf{Workload recordings are generated.} Traffic came from two generators, not real users, and a fault that begins during a very large change of load can be masked by the load-aware model. \emph{Reply:} ``That is why the Hybrid also lets the rules and container facts raise the alarm.''

\item \textbf{We tuned the settings.} Real-system settings were not tuned on a held-out real system, so those numbers are optimistic by an unmeasured margin. The load-model settings were swept on all recordings together. \emph{Reply:} ``We chose the flat region of the sweep, not its maximum, and the synthetic numbers come from a split scored once.''

\item \textbf{A mild fault longer than two minutes is absorbed by the ML.} The rules and container facts keep alarming. \emph{Reply:} ``A trade for not alarming forever on a service that really changed; we measured the alternative (4--20 minutes) and it doubled alarm time with no accuracy gain.''

\item \textbf{Early warning is weak.} Warned before the alarm for 13 of 178 synthetic faults and none of the 62 real ones. \emph{Reply:} ``It is the weakest component and we say so.''

\item \textbf{Auto-restart exists.} The mapper restarts crashed containers by default. \emph{Reply:} ``It had no measurable effect on our benches; switch it off when measuring detection.''

\item \textbf{Not tested on Kubernetes.}

\item \textbf{The ingest endpoint has a built-in token} for local use. \emph{Reply:} ``It must be changed when the server is reachable by others.''
\end{enumerate}

\section{The things that did not work (we reported them)}
\begin{simplebox}
An honest report lists failures. Ours does:

\begin{itemize}
\item The corrected ML engine is \emph{not} more accurate than the first at catching injected faults.
\item Hybrid over Rules was significant in the first version ($p=0.032$); with the corrected engine it is not ($p=0.117$).
\item The random walk is not significantly better than ranking by anomaly alone on our benchmarks.
\item Learned directions do not beat the name rule when names exist.
\item Onset-precedence re-ranking made no difference.
\item The subspace residual (the principled multivariate alternative) is as good at discrimination but raises more false alarms on the load steps and costs more.
\item Early warning does not fire on abrupt real faults.
\item A range guard on the new model added false alarms and was dropped.
\end{itemize}
\end{simplebox}

\section{Claims you should \emph{not} make}
\begin{warnbox}
\begin{itemize}
\item ``Our model is more accurate than the old one at detecting faults.'' (Same accuracy; it is more \emph{robust}.)
\item ``Hybrid is significantly better than rules.'' (Not significant.)
\item ``It predicts failures.'' (It forecasts \emph{who falls next} given a failure; early warning for gradual faults is weak.)
\item ``It works in production.'' (Two benches and a simulator; no production system.)
\item ``It is distributed tracing.'' (It is not.)
\item ``It is autonomous.'' (Its only action is a restart, which can be switched off.)
\item ``Learned propagation works on real cascades.'' (Untested on real cascades.)
\item ``We proved architecture independence.'' (Argued, not proven.)
\end{itemize}
\end{warnbox}

\section{Claims you \emph{can} make with confidence}
\begin{saybox}
\begin{itemize}
\item ``The first ML engine mistook a workload change for a fault: 23--52\% of the time in alarm on recordings with no faults. The corrected engine: 0\%.''
\item ``Container facts are the most valuable single component: 18 points on the real systems.''
\item ``The tool needs no labels, no GPU and no changes to the application; one decision costs about 0.37 ms per service.''
\item ``Learned propagation is better calibrated than fixed edge probabilities: ECE 0.033 against 0.051 on 69 cascades.''
\item ``Learned directions stay correct when names carry no information: 27 of 31 pairs, against none for the name rule.''
\item ``Two measurement bugs inflated the engines in our own benchmark; we found and fixed both.''
\end{itemize}
\end{saybox}
